% $Id: introduction.tex 1784 2012-04-27 23:29:31Z  $
% !TEX root = main.tex



\chapter{Introduction}

Distributed systems replicate data across nodes to allow performance and availability. However, keeping all replicas consistent when dealing with concurrency, delays, and network partitions is a significant challenge. Traditional solutions that enforce strong consistency usually rely on consensus protocols and global coordination, which can be expensive and unstable under failure scenarios. As an alternative solution, \emph{eventual consistency} is proposed. In this case updates may be applied independently at different replicas and reunified later in the process.

\emph{Conflict-Free Replicated Data Types} (CRDTs) were defined to formalize this idea. CRDTs are data structures that allow updating concurrently and can be applied without coordination but still guarantee that all the replicas eventually reach the same state, as long as they receive the same set of updates. This property, known as \emph{Strong Eventual Consistency} (SEC) is achieved through the design of asynchronous update and merge operations which have to satisfy algebraic conditions including commutativity, associativity, idempotence, and monotonicity ~\cite{shapiro2011conflict}.

Designing such data types is not always straightforward. Even small design choices, for instance how to handle removals in sets, can introduce issues that may conduct to divergence or break the intended semantics. Formal verification provides a solution to avoid behavioral errors by proving that a CRDTs satisfy properties like convergence and commutativity. 

Prior work has demonstrated how to specify and verify CRDTs in theorem provers like Isabelle/HOL \cite{zeller2014formal}, establishing rigorous correctness guarantees for individual CRDT implementations. This thesis explores a structured approach by suggesting a more general purpose framework in the Athena proof assistant for specifying and verifying multiples types of CRDTs. The goal is to reduce the effort of building dedicated frameworks to verify each independent implementation of CRDTs by identifying common patterns in their design and behavior, using these patterns in reusable components. Our proposed framework aims to support the definition of state-based CRDTs and helps express their semantics in a modular way. It also provides the structure to reason about properties like convergence and idempotence.

The contribution is not a complete or fully automated system but a set of structures that make it easier to work with CRDTs in Athena. In particular, it is aimed at situations where developers want to experiment with CRDT variants and need a way to express their behaviors precisely and check key correctness properties. While this framework is limited in scope and not intended to cover every possible CRDT design; it offers a modest step toward making formal verification more practical in this domain.

\textbf{Needs a paragraph discussing the validation and results.}
\textbf{Needs a paragraph discussing the thesis outline (structure and what you find in each chapter).}

\endinput